\documentclass[11pt,a4paper]{article}
\usepackage[margin=19mm]{geometry}
\emergencystretch=2em
\usepackage{amsmath,amssymb,booktabs,array}
\usepackage[colorlinks=true,urlcolor=blue,linkcolor=blue]{hyperref}
\setlength{\parindent}{0pt}
\setlength{\parskip}{5pt}
\renewcommand{\arraystretch}{1.12}
\newcommand{\um}{\ensuremath{\,\mu\mathrm{m}}}
\newcommand{\uW}{\ensuremath{\,\mu\mathrm{W}}}
\title{2024 Introduction to VLSI Design\\Homework 3: 4-bit Multiplexer-based Adder}
\author{Reconstructed and electrically verified solution}
\date{Simulation: 13 October 2026}
\begin{document}
\maketitle
\section{Task and final design}
Implement Figure 1 of the assignment using CIC018 TT, $V_{DD}=1.8$ V,
$W\geq0.47\um$, and $L\geq0.18\um$. Preserve the supplied pin order,
input PWL waveforms, $C_{IN}=0$, eight ordered vectors, and 9 ns power window.
Report the worst delay for those vectors, average power, energy, and
$\sum W_iL_i$. No layout or extracted parasitics enter this HW3 result.

The selected design contains \textbf{64 MOSFETs}, all $W=0.47\um$,
$L=0.18\um$. Define
\begin{align}
e_i&=\overline{X_i\oplus Y_i},\\
S_i&=\overline{e_i\oplus C_i}=X_i\oplus Y_i\oplus C_i,\\
C_{i+1}&=e_iY_i+\overline{e_i}C_i.
\end{align}
When $X_i=Y_i$, carry must equal $Y_i$; otherwise carry equals the incoming
$C_i$. Thus the mux equation implements the full-adder carry exactly.
All upper XNORs evaluate in parallel, while carry traverses four serial muxes.
This follows the assignment's serial-parallel diagram; it has no clocked
bit-serial operation or lookahead carry network.

\begin{center}
\setlength{\unitlength}{0.95pt}
\begin{picture}(470,100)
\put(4,43){$C_{IN}$}
\put(30,46){\vector(1,0){20}}
\put(50,28){\framebox(65,36){bit 0}}
\put(115,46){\vector(1,0){40}}
\put(124,52){$C_1$}
\put(155,28){\framebox(65,36){bit 1}}
\put(220,46){\vector(1,0){40}}
\put(229,52){$C_2$}
\put(260,28){\framebox(65,36){bit 2}}
\put(325,46){\vector(1,0){40}}
\put(334,52){$C_3$}
\put(365,28){\framebox(65,36){bit 3}}
\put(430,46){\vector(1,0){20}}
\put(452,43){$C_{OUT}$}
\put(65,83){$X_0,Y_0$}\put(82,80){\vector(0,-1){16}}
\put(170,83){$X_1,Y_1$}\put(187,80){\vector(0,-1){16}}
\put(275,83){$X_2,Y_2$}\put(292,80){\vector(0,-1){16}}
\put(380,83){$X_3,Y_3$}\put(397,80){\vector(0,-1){16}}
\put(82,28){\vector(0,-1){15}}\put(75,0){$S_0$}
\put(187,28){\vector(0,-1){15}}\put(180,0){$S_1$}
\put(292,28){\vector(0,-1){15}}\put(285,0){$S_2$}
\put(397,28){\vector(0,-1){15}}\put(390,0){$S_3$}
\end{picture}
\end{center}

\subsection*{Cell structure and area}
Each bit uses a 4T upper XNOR, a 2T inverter producing $\overline e_i$,
a 2T inverter producing $\overline C_i$, two 2T transmission gates for sum,
and two 2T transmission gates for carry: $4+2+2+4+4=16$ transistors.
\[
A_{\rm gate}=4(16)(0.47)(0.18)=\boxed{5.4144\,\mu\mathrm m^2}.
\]
This is the specified gate-area proxy, not a placed-and-routed footprint.
For comparison, the supplied HW3 88T, $L=0.24\um$ topology has
$A=88(0.47)(0.24)=9.9264\,\mu\mathrm m^2$.

\newpage
\section{Transistor implementation}
The compact upper XNOR can have a threshold-degraded \emph{internal} high
when $X=Y=1$. The lower XNOR uses $C_i/\overline C_i$ as its \emph{data}
and $e_i/\overline e_i$ as its \emph{select}; it does not pass the degraded
$e_i$ directly to $S_i$. The complementary CMOS inverter and PMOS/NMOS
transmission-gate pair allow the external outputs to approach the rails.
The actual TT electrical results, rather than an ideal Boolean assumption,
establish the output levels below. This design has not been characterized
for other corners, supplies, loads, or mismatch.

The following complete DUT source is included separately as
\path{four_bit_multbased_adder.sp}. The testbench provides global rails
and loads the licensed model library; the model is not part of the repository.

{\footnotesize\begin{verbatim}
* Figure 1 serial-parallel multiplexer-based adder, 64 MOSFETs.
* Dimensions in micrometres; global rails vdd/vss provided by testbench.
.param WN=0.47u WP=0.47u WT=0.47u LG=0.18u
.subckt inv a z
MP z a vdd vdd P_18 W=WP L=LG
MN z a vss vss N_18 W=WN L=LG
.ends inv
.subckt tg a z s sb
MP z sb a vdd P_18 W=WT L=LG
MN z s a vss N_18 W=WT L=LG
.ends tg
.subckt mux_fa x y ci s co
* Compact 4T upper XNOR: e can be threshold-degraded internally.
MPU0 ep x vdd vdd P_18 W=WP L=LG
MPU1 e y ep vdd P_18 W=WP L=LG
MNU0 e y x vss N_18 W=WN L=LG
MNU1 e x y vss N_18 W=WN L=LG
XI2 e en inv
XI3 ci cib inv
* Lower XNOR uses ci/cib as DATA; e/en as SELECT.
* This choice restores the external sum even when compact e is degraded.
XT2 ci s e en tg
XT3 cib s en e tg
* Carry: co = e ? y : ci. Shared complementary select.
XT4 y co e en tg
XT5 ci co en e tg
.ends mux_fa
.subckt four_bit_multbased_adder cin x0 x1 x2 x3 y0 y1 y2 y3 s0 s1 s2 s3 cout
XF0 x0 y0 cin s0 c1 mux_fa
XF1 x1 y1 c1 s1 c2 mux_fa
XF2 x2 y2 c2 s2 c3 mux_fa
XF3 x3 y3 c3 s3 cout mux_fa
.ends four_bit_multbased_adder
\end{verbatim}}

\newpage
\section{Measurement definitions and required-vector results}
All results use PrimeSim HSPICE S-2021.09 win64, CIC018 TT, 25$^\circ$C,
and 1.8 V, with no added output capacitors. Power includes the complete
0--9 ns interval, including the initial zero-vector interval and glitches.
\[
P_{\rm total}=\operatorname{AVG}(\texttt{POWER}),\quad
P_{VDD}=\operatorname{AVG}[-V(vdd)I(vvdd)],\quad E_9=P_{\rm total}(9\,\mathrm{ns}).
\]
HSPICE's total circuit POWER is the negative sum of the powers of all
independent sources. Input PWL drivers can supply energy through TG data
paths, so $P_{\rm total}$ differs from $P_{VDD}$. The assignment's
\texttt{.meas tran pwr avg power} uses total power; substituting the smaller
VDD-only value would change the grading metric.

\begin{center}\small
\begin{tabular}{lrrr}\toprule
Measurement setting & $P_{\rm total}$ ($\mu$W) & $P_{VDD}$ ($\mu$W) & $E_9$ (fJ)\\\midrule
Original \texttt{.tran 0.1ns 9n} & 69.002216 & 37.998479 & 621.019943\\
Refined: 1 ps print / maximum step & 69.558498 & 37.432840 & 626.026482\\\bottomrule
\end{tabular}\end{center}
The original deck leaves DELMAX automatic. Refined measurements use
\texttt{.tran 1p 9n} and \texttt{.option delmax=1p}; the 0.806\% total-energy
difference is reported explicitly. Candidate comparisons and delays use
the same refined setting. No PWL input time or voltage is changed.

\begin{center}\small
\begin{tabular}{rrrrrl}\toprule
Vector & $X$ & $Y$ & Expected $\{C_{OUT},S_3,S_2,S_1,S_0\}$ & Measured & Result\\\midrule
1 & 1 & 5 & 00110 & 00110 & pass \\
2 & 15 & 13 & 11100 & 11100 & pass \\
3 & 11 & 9 & 10100 & 10100 & pass \\
4 & 4 & 8 & 01100 & 01100 & pass \\
5 & 3 & 7 & 01010 & 01010 & pass \\
6 & 12 & 14 & 11010 & 11010 & pass \\
7 & 2 & 15 & 10001 & 10001 & pass \\
8 & 6 & 10 & 10000 & 10000 & pass \\
\bottomrule\end{tabular}\end{center}
At 0.8 ns after each vector endpoint, all eight vectors pass
$V_L\leq0.2$ V and $V_H\geq1.6$ V. The maximum sampled low is
0.00396849 V and minimum high is 1.7910275 V. These are actual finite-time
voltage readings, not an assertion of mathematically exact 0/1.8 V outputs.

\subsection*{Causal delay and glitches}
For vector $k$, all changed input ramps run from $k-0.1$ to $k$ ns;
their 50\% crossings occur at $k-0.05$ ns. For every output whose expected
bit changes, use its \emph{last} 0.9 V crossing in that vector interval:
\[
t_{pd}=t_{\rm out,last\ 50\%}-(k-0.05\,\mathrm{ns}).
\]
The maximum is \textbf{152.910320 ps}, from $(2,15)\rightarrow(6,10)$,
$S_0$ falling: $t_{in}=7.950$ ns and $t_{out}=8.102910320$ ns.
This is the worst observed delay for the required sequence, not a global
maximum over all possible input transitions. The output data contain 12
extra 50\% crossings from glitches. The conservative 1 ps sampled bound
for all outputs to settle in the correct 0.2/1.6 V bands is 394 ps relative
to the common input 50\% event; it is a different metric from $t_{pd}$.

\newpage
\section{Delay evidence and sizing comparison}
\begin{center}\small
\begin{tabular}{rllrrr}\toprule
Vector & Output & Edge & Input 50\% (ns) & Last output 50\% (ns) & Delay (ps)\\\midrule
1 & S1 & rise & 0.950 & 0.985495 & 35.495 \\
1 & S2 & rise & 0.950 & 1.041877 & 91.877 \\
2 & S1 & fall & 1.950 & 2.053628 & 103.628 \\
2 & S3 & rise & 1.950 & 2.088414 & 138.414 \\
2 & COUT & rise & 1.950 & 1.956704 & 6.704 \\
3 & S3 & fall & 2.950 & 3.020838 & 70.838 \\
4 & S3 & rise & 3.950 & 4.044057 & 94.057 \\
4 & COUT & fall & 3.950 & 4.067813 & 117.813 \\
5 & S1 & rise & 4.950 & 4.981098 & 31.098 \\
5 & S2 & fall & 4.950 & 5.030657 & 80.657 \\
6 & COUT & rise & 5.950 & 5.956614 & 6.614 \\
7 & S0 & rise & 6.950 & 7.040801 & 90.801 \\
7 & S1 & fall & 6.950 & 7.079615 & 129.615 \\
7 & S3 & fall & 6.950 & 7.057709 & 107.709 \\
8 & S0 & fall & 7.950 & 8.102910 & 152.910 \\
\bottomrule\end{tabular}\end{center}

The original 88T design uses independent select inverters in its lower
XNOR and carry mux. Sharing the complement reduces this to 80T. An 8T
full-swing upper XNOR is retained in the 80T candidates. A second family
uses the compact 4T upper XNOR from the supplied 2024 HW4 schematic idea,
reducing the total to 64T. All six candidates pass the eight required
vectors electrically under the same refined measurement conditions.

\begin{center}\small
\begin{tabular}{lrrrr}\toprule
Candidate & $\sum WL$ ($\mu$m$^2$) & $E_9$ (fJ) & $t_{pd}$ (ps) & $P t_{pd}$ (fJ)\\\midrule
88T reference-equivalent & 9.9264 & 1005.327 & 247.239 & 27.617 \\
80T minimum & 6.7680 & 824.288 & 137.009 & 12.548 \\
80T inverter P=0.94 & 8.1216 & 947.745 & 133.867 & 14.097 \\
80T TG W=0.70 & 8.7552 & 1055.585 & 161.400 & 18.930 \\
64T minimum (selected) & 5.4144 & 626.026 & 152.910 & 10.636 \\
64T P=0.70 & 6.0768 & 676.239 & 149.882 & 11.262 \\
\bottomrule\end{tabular}\end{center}
The selected minimum-size 64T circuit has the smallest area, energy,
and power--delay product in this finite candidate set. Relative to the
reference-equivalent 88T circuit, energy decreases by 37.73\%, area by
45.45\%, and measured delay by 38.15\%. The minimum-size 80T circuit is
faster than the 64T circuit, but consumes more energy and area and has a
higher power--delay product. Increasing PMOS or TG sizes does not improve
the final combined choice. This comparison does not establish a global
optimum over all transistor networks or continuous sizes.

\newpage
\section{Electrical exhaustive verification and reproduction}
In addition to the supplied sequence, the selected transistor circuit was
simulated for all $16\times16\times2=\textbf{512}$ combinations of $X$, $Y$,
and $C_{IN}$. Each vector occupies 2 ns, with 0.1 ns input ramps;
voltages are sampled 1.8 ns after the endpoint. All \textbf{512/512} cases
produce the arithmetic result $X+Y+C_{IN}$ and satisfy $V_L\leq0.2$ V,
$V_H\geq1.6$ V on all five external outputs. Across those samples,
$\max V_L=0.002939216$ V and $\min V_H=1.79601867$ V.
This is transistor-level transient sampling of every settled input
combination, not just Boolean software evaluation. It does not exhaust
all $512^2$ transitions or prove global worst delay.

\texttt{simulate.py} generates six candidate decks, runs licensed HSPICE,
exports measured values, selects the 64T circuit, verifies the 512 electrical
cases, and runs the original instructor testbench separately. Set
\texttt{CIC018\_MODEL} or \texttt{--model} to the licensed library and
\texttt{HSPICE} or \texttt{--hspice} to the simulator executable.
\begin{verbatim}
python simulate.py --model D:/cic018.l --work .simulation --results results
\end{verbatim}
The default Windows HSPICE installation is supplied in the script and can
be overridden. All transistor netlists use the required
\texttt{P\_18}/\texttt{N\_18} models and the exact assignment subcircuit pin order.
The standalone original testbench changes only the model path from the
instructor's Linux path to \texttt{D:/cic018.l}.

\begin{center}\small
\begin{tabular}{p{0.35\linewidth}p{0.57\linewidth}}\toprule
Output file & Evidence or purpose\\\midrule
\texttt{results/candidates.csv} & Dimensions, device count, area, total and supply power, energy, delay, objective products.\\
\texttt{results/vectors.csv} & Eight required vector results with five actual output voltages.\\
\texttt{results/delays.csv} & Causal input and output crossing times for every expected-changing output.\\
\texttt{results/settling.csv} & Per-vector output crossing counts and sampled band-settling bounds.\\
\texttt{results/exhaustive\_512.csv} & All input combinations, arithmetic result, actual voltages, pass flags.\\
\texttt{results/waveform.csv} & 1 ps sampled output voltages and signed VDD current; GUI screenshots are unnecessary.\\
\texttt{results/summary.json} & Conditions, model SHA256 identity, limits and result summary.\\
\texttt{results/instructor\_exact.json} & Original 100 ps print-step / automatic DELMAX power result.\\\bottomrule
\end{tabular}\end{center}

\newpage
\section*{Sources reviewed and integration decisions}
\begin{enumerate}
\item NCU, \emph{EE3032 Introduction to VLSI Design,
Homework 3}, Fall 2024; supplied Figure 1, original testbench and explanatory
notes determine the topology, stimuli, minimum width, and grading window.
\item User-supplied 2024 HW3 reference by Chih-Hua Chang, archive containing
\path{four_bit_multbased_adder.sp} and report. Its 88T TG design is
reconstructed for equal-condition baseline comparison; misleading XOR
comments in that source do not change its actual XNOR wiring.
\item User-supplied 2024 HW4 reference by Chih-Hua Chang, schematic
\path{four_bit_multbased_adder.sp}. Its compact upper-XNOR idea is
integrated after new electrical validation; neither its old measurements
nor extracted layout parasitics are reused as HW3 results.
\item Public source: \href{https://github.com/dOwOb1212/NCU_VLSI_Design_2023_HW/blob/main/hw1/CLA_cmos.sp}{dOwOb1212, NCU\_VLSI\_Design\_2023\_HW, hw1/CLA\_cmos.sp}.
Despite its CLA label, carry is passed serially through muxes. The source
is 8-bit, uses widths below the 2024 minimum, and has different stimuli;
it is an architectural cross-check, not a valid drop-in 2024 submission.
\item Local Synopsys S-2021.09 documentation: \emph{Selecting Simulation
Output Parameters}, \texttt{.OPTION DELMAX}, and \texttt{.OPTION LIS\_NEW}.
These define total POWER, internal time-step control, and machine-readable
\texttt{.printtr0} output.
\end{enumerate}
All numerical results above were newly simulated on 13 October 2026.
The public repository contains derived statistics and source only; the
licensed process library, student original reports, and raw license logs
remain local.
\end{document}
